\uuid{wTlx}
\titre{Primitives usuelles}
\auteur{hanani ; mzari}
\organisation{exo7}
\licence{CC BY-NC-SA 4.0}
\datecreate{2020-01-20}
\niveau{L1}
\module{Analyse}
\chapitre{Calcul d'intégrales}
\sousChapitre{Primitives diverses}
\difficulte{1}
\theme{Primitives, primitives usuelles}
\exo7qcmid{lille-2.primitives.001}

\contenu{
\begin{question}
\qclassification{127.11}
\qtags{motcle=primitives usuelles}
Parmi les affirmations suivantes, cocher celles qui sont vraies :
\begin{answers}
\good{Si $u$ est une fonction dérivable, strictement positive, alors $\sqrt{u}$ est une primitive de $\displaystyle \frac{u'}{2\sqrt{u}}$.}
\good{Si $u$ est une fonction dérivable, alors $\arctan (u)$ est une primitive de $\displaystyle \frac{u'}{1+u^2}$.}
\bad{Si $u$ est une fonction dérivable, alors $\mathrm{e}^u$ est une primitive de $\displaystyle \mathrm{e}^u$.}
\bad{Si $u$ est une fonction dérivable et ne s'annulant pas, alors $\ln (u)$ est une primitive de $\displaystyle \frac{u'}{u}$.}
\end{answers}
\begin{explanations}
On a : $\displaystyle (\sqrt{u})'=\frac{u'}{2\sqrt{u}}$, $\displaystyle (\arctan (u))'=\frac{u'}{1+u^2}$ et $\left(\mathrm{e}^u\right)'=u'\mathrm{e}^u$. Donc $\displaystyle 2\sqrt{u}$ est une primitive de $\displaystyle \frac{u'}{\sqrt{u}}$, $\displaystyle \arctan (u)$ est une primitive de $\displaystyle \frac{u'}{1+u^2}$ et $\mathrm{e}^u$ n'est pas une primitive de $\displaystyle \mathrm{e}^u$. Enfin, $\ln (u)$ peut ne pas être définie. Il suffit de prendre $u=-1-x^2$, par exemple.
\end{explanations}
\end{question}
}
